\documentclass[12pt,a4paper]{article}
\usepackage[utf8]{inputenc}\usepackage[T1]{fontenc}\usepackage[provide=*,french]{babel}
\usepackage{lmodern,microtype,longtable,booktabs,array,enumitem,fancyhdr,needspace}
\usepackage[margin=2.1cm,headheight=16pt]{geometry}
\usepackage[hidelinks,breaklinks=true]{hyperref}
\setlength{\parindent}{0pt}\setlength{\parskip}{6pt}\setlength{\emergencystretch}{3em}
\setlist[itemize]{leftmargin=1.3em,itemsep=3pt}
\pagestyle{fancy}\fancyhf{}\fancyhead[L]{\small Master Finance verte — Module 6}\fancyhead[R]{\small Badr Aboufarasse}\fancyfoot[C]{\thepage}
\begin{document}
\begin{center}{\small MASTER FINANCE VERTE -- MODULE 6}\par\vspace{8pt}{\Large\bfseries Séance 2 -- Plan financier et trajectoire de transition}\par\vspace{6pt}Badr Aboufarasse\par{\small Cours et applications corrigées -- 25 octobre 2026}\end{center}
\begin{quote}\itshape Combien financer, à quelle date et pour quels résultats ?\end{quote}
\textbf{Données du cas fictives.} Montants en M MAD ; taux et seuils pédagogiques.

\subsection*{Objectifs}
\begin{itemize}\item Construire un plan financier reproductible sur cinq ans.\item Relier les trajectoires physiques aux flux et aux dates de mise en service.\item Distinguer besoin annuel et risque de liquidité en cours d’année.\end{itemize}
\section*{1. Le plan comme système d’hypothèses}
Un plan financier n’est pas une extrapolation mécanique du résultat. Il décrit des décisions, leur calendrier et leurs effets. On commence par les facteurs économiques : volumes, prix, coûts, délais clients, stocks et investissements. Les états prévisionnels en découlent. Chaque hypothèse doit posséder une unité, une date, une justification et un responsable de révision. [BD] [SS7]


Dans le cas, le CA progresse de 5 \% par an ; les charges variables sont 60 \% du CA ; les charges fixes de 24 augmentent de 2 \% par an ; le BFR représente 20 \% du CA. Cette dernière convention est une simplification de délais et de stocks stables. Elle ne prouve pas que tous les postes du BFR évoluent toujours proportionnellement aux ventes.


Le modèle est nominal en MAD. Les hausses de 5 \% et 2 \% sont déjà nominales ; on n’ajoute pas une inflation générale une seconde fois. Les hypothèses d’impôt et de financement sont pédagogiques. Une modélisation réelle séparerait prix, volumes, devises et différents taux d’inflation.

\section*{2. Le programme commun et son calendrier}

\Needspace{215pt}
\begin{longtable}{p{3.67cm}p{3.67cm}p{3.67cm}p{3.67cm}}
\toprule
Investissement & Décaissement & Gain annuel d’EBITDA & Début du gain \\
\midrule\endhead
Maintenance & 4 chaque année & Préserve la capacité & Déjà dans la référence \\
Efficacité E & 6 en A1 & 1,8 & A2 \\
Solaire S détenu & 10 en A1 & 2,0 net de maintenance spécifique & A2 \\
Eau W & 4 en A2 & 0,7 & A3 \\
\bottomrule
\end{longtable}

CAPEX total : 20 en A1, 8 en A2, puis 4 par an. Amortissement historique maintenu à 4 ; E et S ajoutent 0,6 et 1 à partir de A2 ; W ajoute 0,4 à partir de A3. Ces durées de dix ans et dates sont des conventions explicites du cas. Les gains sont constants en montant nominal pour isoler le mécanisme.


Les flux annuels sont agrégés et les intérêts calculés sur les encours d’ouverture, nouvelle dette comprise à partir de A1. Le remboursement de principal se fait en fin d’année. Le modèle ne mesure pas les pics infra-annuels : le zoom mensuel doit être traité séparément.


La maintenance historique de 4 ne s’ajoute pas une seconde fois au gain net solaire de 2. Celui-ci correspond à une dépense énergétique évitée de 2,4 moins 0,4 de maintenance spécifique. La référence sans projets ne comprend pas les économies E, S et W.

\section*{3. Construire le compte de résultat prévisionnel}

\begin{quote}\small CA(t) = 100 $\times$ 1,05\textasciicircum{}t\newline EBITDA(t) = 40 \% $\times$ CA(t) - 24 $\times$ 1,02\textasciicircum{}t + économies actives\newline EBIT = EBITDA - amortissements\newline Impôt cash = max[0 ; 25 \% $\times$ (EBIT - intérêts)]\end{quote}

L’impôt est calculé après intérêts dans ce plan de trésorerie. Aucun crédit d’impôt immédiat n’est reconnu en cas de perte ; les reports fiscaux sont exclus. Ce choix évite de transformer une perte fiscale en encaissement artificiel.


En A1 : CA = 105 ; charges variables = 63 ; charges fixes = 24,48 ; EBITDA = \textbf{17,52}. Aucun gain de transition n’est encore actif. L’EBIT est 13,52. Avec une dette existante de 22 à 6 \% et un nouveau prêt de 14 à 6 \%, les intérêts sont 1,32 + 0,84 = \textbf{2,16}. L’impôt vaut 2,84 ; le résultat net \textbf{8,52}.


Le taux historique implicite ne devient pas automatiquement le taux futur : les intérêts N de 1,8 sont une donnée historique. Le taux contractuel futur de 6 \% est une hypothèse distincte. Il faut expliquer cette séparation pour éviter une incohérence apparente.

\section*{4. Du résultat aux ressources à mobiliser}

\begin{quote}\small CFO = résultat net + amortissements - $\Delta$BFR\newline Cash final = cash initial + CFO - CAPEX + dette tirée + apport - principal remboursé - dividendes\end{quote}

La politique de référence retient aucun dividende de A1 à A5, aucun apport et un nouveau prêt de 14 tiré au début de A1. La dette existante rembourse 2 par an ; la nouvelle 2,8 par an pendant cinq ans. En A1, $\Delta$BFR = 21 - 20 = 1 et CFO = 8,52 + 4 - 1 = \textbf{11,52}. La trésorerie de clôture est 4 + 11,52 - 20 + 14 - 4,8 = \textbf{4,72}.


Cette clôture dépasse le coussin de 3 retenu par l’entreprise. Ce coussin est une politique pédagogique, pas un seuil bancaire universel. Si le nouveau prêt était absent, il faudrait recalculer les intérêts et l’impôt : on ne peut pas seulement retrancher 14 du cash final, car les frais financiers changent également.


Le financement doit être lié au modèle : un emprunt peut résoudre une impasse initiale puis créer une échéance trop lourde. Les intérêts dépendent ici du solde d’ouverture pour éviter une circularité artificielle entre tirage de dette, intérêts et déficit de financement.

\section*{5. Résultats de référence et contrôles}

\Needspace{250pt}
\begin{longtable}{p{2.84cm}p{2.84cm}p{2.84cm}p{2.84cm}p{2.84cm}}
\toprule
Année & CA & EBITDA & Cash final & DSCR \\
\midrule\endhead
A1 & 105.000 & 17.520 & 4.720 & 1.391 \\
A2 & 110.250 & 22.930 & 8.064 & 2.101 \\
A3 & 115.763 & 25.336 & 17.475 & 2.474 \\
A4 & 121.551 & 27.142 & 28.402 & 2.792 \\
A5 & 127.628 & 29.053 & 40.921 & 3.155 \\
\bottomrule
\end{longtable}

Les fichiers CSV téléchargeables donnent tous les postes et les résultats non arrondis. Les économies E et S commencent en A2 ; W commence en A3. La trésorerie s’accumule ensuite parce que le cas ne prévoit pas de nouvelle croissance, de dividende ou d’acquisition. Ce solde ne constitue pas une recommandation de thésaurisation.


\textbf{Contrôles :} dette finale = dette initiale + tirages - remboursements ; cash final = cash initial + flux nets ; gains nuls avant mise en service ; CAPEX de maintenance présent chaque année ; absence de double retrait des intérêts ; hypothèses de BFR identiques dans les résultats et le cash.


Le plan vise la traçabilité, pas une précision illusoire. Il faut pouvoir modifier une hypothèse et expliquer le sens du résultat. Les valeurs arrondies affichées ne doivent pas servir de base aux calculs suivants.

\section*{6. Le risque de calendrier et les informations de transition}
Une clôture annuelle positive peut masquer une rupture de liquidité. Les salaires sont mensuels, un acompte fournisseur peut être exigé avant la livraison et une banque peut décaisser après constat d’avancement. L’échéancier d’investissement et les conditions de mobilisation du financement doivent donc être rapprochés. Une ligne non confirmée n’est pas une réserve certaine.


La trajectoire de transition doit aussi expliquer quelles dépenses rendent possible quel résultat physique. Un budget « vert » sans dates de mise en service ne permet pas de relier objectifs et moyens. Les ressources IFRS sur la transition aident à articuler stratégie, effets financiers et incertitudes. Elles n’imposent pas par elles-mêmes à toute entreprise de créer un plan de transition. [IFRS]


Dans le dossier, indiquer pour E, S et W : dépense, mise en service, effet opérationnel, indicateur physique et condition de réalisation. La réduction d’émissions sera étudiée séparément de la réduction de facture : un prix électrique n’est pas un facteur d’émission.

\section*{7. Trois applications corrigées}

\subsection*{Application 1 --- Une croissance qui consomme du cash}
Avec BFR/CA constant à 20 \%, calculer le BFR en A1 et A2 et les utilisations de cash correspondantes.


\paragraph{Correction expliquée}
CA A1 = 105 ; A2 = 110,25. BFR = 21 puis 22,05 ; variations = \textbf{1 puis 1,05}. Le BFR de 22,05 n’est pas intégralement décaissé chaque année : seule sa variation est un flux de l’année.



\subsection*{Application 2 --- Pic de liquidité}
Au premier mois, le cash initial est 4, les encaissements 7, les paiements d’exploitation et impôt 9, les paiements d’investissement 5 et le tirage de dette 4. Il n’y a pas d’autre flux. Quel est le besoin pour préserver un coussin de 3 ?


\paragraph{Correction expliquée}
Cash final = 4 + 7 - 9 - 5 + 4 = \textbf{1}. Besoin complémentaire = \textbf{2}. Le cas mensuel illustre un profil de flux ; il n’est pas une ventilation exacte de toute la prévision annuelle. Proposer un tirage anticipé confirmé, un échéancier fournisseur négocié ou un décalage du programme, puis chiffrer leur coût.



\subsection*{Application 3 --- Retard des gains}
E et S entrent finalement en service en A3, avec décalage des amortissements également à A3. Tous les autres paramètres et le CAPEX sont inchangés. Quel est l’effet sur le CFO A2, si l’entreprise reste bénéficiaire ?


\paragraph{Correction expliquée}
La perte d’EBITDA est 3,8 et l’amortissement baisse de 1,6. L’EBIT et le résultat imposable baissent de 2,2 ; l’impôt baisse de 0,55. CFO diminue donc de \textbf{3,25} (= 3,8 - 0,55). Retirer seulement 3,8 du cash ignorerait l’impôt ; utiliser un bouclier fiscal d’amortissement malgré le report serait également erroné.

\section*{8. Cas de synthèse et transition}
\textbf{Atelier, 30 minutes.} Reproduire A1 puis A2 dans un tableur à partir du dossier commun. Séparer cellules d’hypothèses et formules. Rédiger une note expliquant si le prêt de 14 suffit à la clôture annuelle, puis pourquoi cela ne démontre pas l’absence d’impasse mensuelle.



\paragraph{Éléments de solution}
En A2 : charges fixes 24 $\times$ 1,02² = 24,9696 ; EBITDA = 44,1 - 24,9696 + 3,8 = \textbf{22,9304} ; amortissements 5,6 ; intérêts 1,872 ; impôt 3,8646 ; résultat net 11,5938 ; CFO 16,1438. Le cash final est 4,72 + 16,1438 - 8 - 4,8 = \textbf{8,0638}. Le prêt couvre le plan aux clôtures de référence. Un échéancier mensuel reste nécessaire pour prouver la disponibilité des fonds au moment des paiements.


\textbf{À retenir :} une trajectoire devient finançable lorsque les moyens sont disponibles aux dates requises et que les engagements restent supportables. La séance 3 compare les structures de financement au lieu de considérer le prêt comme une simple cellule d’ajustement.

\section*{Références mobilisées}
\small
\begin{itemize}
\item \textbf{[BD] Berk, J. et DeMarzo, P..} Corporate Finance, 6e édition, Pearson. Chapitres 2, 12, 14--18, 22, 24--27 et 30--31 selon la séance. \href{https://www.pearson.com/en-us/subject-catalog/p/corporate-finance/P200000009791/9780137844906}{Source officielle}.
\item \textbf{[SS7] Schoenmaker, D. et Schramade, W. (2023).} Capital Budgeting. Chap. 7, p. 173--204. \href{https://doi.org/10.1007/978-3-031-35009-2_7}{Source officielle}.
\item \textbf{[IFRS] IFRS Foundation (2025).} Disclosing information about an entity’s climate-related transition, including information about transition plans, in accordance with IFRS S2. \href{https://www.ifrs.org/content/dam/ifrs/supporting-implementation/ifrs-s2/transition-plan-disclosure-s2.pdf}{Source officielle}.
\end{itemize}\normalsize
\end{document}
